\chapter{麦克斯韦方程组应用}

\section{竞赛题}

\begin{competition}
同轴圆形平板电容器两极板半径为 $R$，极板间距远小于 $R$，忽略边缘效应。若极板间电场可写为
\[
\vb{E}(t)=E_0\sin \omega t\,\vu{z},
\]
求：
\begin{enumerate}
  \item 半径 $r<R$ 的圆周所围面积上的位移电流 $I_d(r,t)$；
  \item 半径 $r$ 处环形磁场 $\vb{B}(r,t)$ 的大小与方向。
\end{enumerate}
\begin{solution}
位移电流定义为
\[
I_d=\epsilon_0\dv{}{t}\Phi_E,
\qquad
\Phi_E=\int \vb{E}\cdot \dd{\vb{S}}.
\]
对半径 $r<R$ 的圆面，电场近似均匀，故
\[
\Phi_E(r,t)=E_0\sin \omega t\cdot \pi r^2.
\]
于是
\[
I_d(r,t)=\epsilon_0\dv{}{t}\qty(E_0\pi r^2\sin \omega t)
=\epsilon_0\pi r^2 E_0\omega \cos \omega t.
\]
再由安培--麦克斯韦方程
\[
\oint \vb{B}\cdot \dd{\vb{l}}=\mu_0 I_d,
\]
利用轴对称性，在半径 $r$ 的圆周上 $B$ 为常量且方向沿周向，故
\[
B(2\pi r)=\mu_0\epsilon_0\pi r^2 E_0\omega \cos \omega t.
\]
解得
\[
B(r,t)=\frac{\mu_0\epsilon_0\omega E_0}{2}r\cos \omega t.
\]
方向由右手定则确定：若 $\pdv*{E}{t}>0$，则 $\vb{B}$ 沿逆时针方向（从 $+z$ 方向看）。
\end{solution}
\end{competition}

\begin{competition}
真空中一列平面电磁波沿 $+x$ 方向传播，电场为
\[
\vb{E}(x,t)=E_0\cos(kx-\omega t)\,\vu{y}.
\]
求：
\begin{enumerate}
  \item 磁场 $\vb{B}(x,t)$；
  \item 瞬时坡印廷矢量 $\vb{S}$；
  \item 一个周期内的平均能流密度 $\avg{S}$。
\end{enumerate}
\begin{solution}
在真空平面波中，$\vb{E}$、$\vb{B}$ 与传播方向两两垂直，且
\[
B_0=\frac{E_0}{c}.
\]
因此
\[
\vb{B}(x,t)=\frac{E_0}{c}\cos(kx-\omega t)\,\vu{z}.
\]
坡印廷矢量为
\[
\vb{S}=\frac{1}{\mu_0}\vb{E}\times\vb{B}.
\]
代入得
\[
\vb{S}(x,t)=\frac{1}{\mu_0}\frac{E_0^2}{c}\cos^2(kx-\omega t)\,\vu{x}
=\epsilon_0 c E_0^2\cos^2(kx-\omega t)\,\vu{x}.
\]
因为一个周期内
\[
\avg{\cos^2(kx-\omega t)}=\frac{1}{2},
\]
所以平均能流密度为
\[
\avg{\vb{S}}=\frac{1}{2}\epsilon_0 c E_0^2\,\vu{x}.
\]
其大小为
\[
\avg{S}=\frac{1}{2}\epsilon_0 c E_0^2.
\]
\end{solution}
\end{competition}

\section{大学物理题}

\begin{university}
真空中一列单色平面电磁波写成
\[
\vb{E}=E_0\cos(kx-\omega t)\,\vu{y},
\qquad
\vb{B}=B_0\cos(kx-\omega t)\,\vu{z}.
\]
利用麦克斯韦方程证明 $E_0=cB_0$，并给出 $k$ 与 $\omega$ 的关系。
\begin{solution}
由法拉第电磁感应定律
\[
\curl \vb{E}=-\pdv{\vb{B}}{t}.
\]
因为 $\vb{E}$ 只有 $y$ 分量，且只依赖于 $x,t$，故
\[
\curl \vb{E}=\pdv{E_y}{x}\vu{z}
=-kE_0\sin(kx-\omega t)\,\vu{z}.
\]
另一方面
\[
-\pdv{\vb{B}}{t}=-\omega B_0\sin(kx-\omega t)\,\vu{z}.
\]
比较系数得
\[
kE_0=\omega B_0.
\]
再由安培--麦克斯韦方程（真空中）
\[
\curl \vb{B}=\mu_0\epsilon_0\pdv{\vb{E}}{t},
\]
可得
\[
kB_0=\mu_0\epsilon_0\omega E_0.
\]
将两式相除得到
\[
\frac{E_0}{B_0}=\frac{\omega}{k}.
\]
而真空波速满足
\[
\frac{\omega}{k}=\frac{1}{\sqrt{\mu_0\epsilon_0}}=c,
\]
故
\[
E_0=cB_0,
\qquad
\omega=ck.
\]
\end{solution}
\end{university}

\begin{university}
某激光束在真空中的平均强度为 $I=5.0\times 10^3\,\mathrm{W/m^2}$，垂直照射在面积 $A=2.0\times 10^{-4}\,\mathrm{m^2}$ 的探测面上。求穿过该面的平均功率，并由坡印廷矢量解释其物理意义。
\begin{solution}
平均强度就是平均坡印廷矢量的模：
\[
I=\avg{S}=\abs{\avg{\vb{S}}}.
\]
因此穿过面积 $A$ 的平均功率为
\[
P=IA=(5.0\times 10^3)(2.0\times 10^{-4})=1.0\,\mathrm{W}.
\]
坡印廷矢量定义为
\[
\vb{S}=\frac{1}{\mu_0}\vb{E}\times\vb{B},
\]
其方向表示电磁能量传播方向，其大小表示单位时间穿过单位面积的电磁能量。故本题中平均每秒有 $1.0\,\mathrm{J}$ 的电磁能通过探测面。
\end{solution}
\end{university}

\begin{university}
真空中平面电磁波的平均强度为 $I$。证明该波的平均动量密度 $\avg{g}$ 与平均辐射压 $p$ 满足：
\[
\avg{g}=\frac{I}{c^2},
\qquad
p=\frac{I}{c}\quad\text{(完全吸收)},
\qquad
p=\frac{2I}{c}\quad\text{(完全反射)}.
\]
\begin{solution}
电磁场动量密度为
\[
\vb{g}=\frac{\vb{S}}{c^2}.
\]
取时间平均即得
\[
\avg{g}=\frac{\avg{S}}{c^2}=\frac{I}{c^2}.
\]
若波垂直入射到表面，单位时间单位面积入射的能量为 $I$，因此单位时间单位面积传递的动量为
\[
\frac{I}{c}.
\]
这正是辐射压，所以完全吸收时
\[
p=\frac{I}{c}.
\]
若完全反射，则法向动量改变量加倍，故
\[
p=\frac{2I}{c}.
\]
这表明辐射压本质上是电磁波所携带动量通量对物体表面的作用。
\end{solution}
\end{university}

\begin{university}
写出电磁场能量守恒的微分形式，并说明每一项的物理意义。进一步对真空情形下的平面波说明能量在电场与磁场之间如何分配。
\begin{solution}
电磁场能量守恒的微分形式即坡印廷定理：
\[
\pdv{u}{t}+\div \vb{S}+\vb{J}\cdot\vb{E}=0,
\]
其中
\[
u=\frac{1}{2}\qty(\epsilon_0 E^2+\frac{B^2}{\mu_0}),
\qquad
\vb{S}=\frac{1}{\mu_0}\vb{E}\times\vb{B}.
\]
各项意义如下：
\begin{itemize}
  \item $\pdv*{u}{t}$：单位体积内电磁能的变化率；
  \item $\div \vb{S}$：电磁能流出该体积的净速率；
  \item $\vb{J}\cdot\vb{E}$：电场对带电粒子做功的功率密度。
\end{itemize}
在真空平面波中，$\vb{J}=0$，故有
\[
\pdv{u}{t}+\div \vb{S}=0.
\]
又因 $E=cB$，可得
\[
\frac{B^2}{\mu_0}=\epsilon_0E^2,
\]
所以电场能密度与磁场能密度相等：
\[
u_E=\frac{1}{2}\epsilon_0E^2,
\qquad
u_B=\frac{1}{2\mu_0}B^2=\frac{1}{2}\epsilon_0E^2.
\]
因此总能量一半储存在电场中，一半储存在磁场中，并由坡印廷矢量沿传播方向输运。
\end{solution}
\end{university}

\section{综合题}

\begin{problem}
从真空中的麦克斯韦方程组出发，完整推导电场与磁场满足的波动方程：
\[
\nabla^2\vb{E}-\mu_0\epsilon_0\pdv[2]{\vb{E}}{t}=0,
\qquad
\nabla^2\vb{B}-\mu_0\epsilon_0\pdv[2]{\vb{B}}{t}=0.
\]
要求写出每一步所用到的矢量恒等式。
\begin{solution}
真空中麦克斯韦方程组为
\[
\div \vb{E}=0,
\qquad
\div \vb{B}=0,
\qquad
\curl \vb{E}=-\pdv{\vb{B}}{t},
\qquad
\curl \vb{B}=\mu_0\epsilon_0\pdv{\vb{E}}{t}.
\]
先推导 $\vb{E}$ 的波动方程。对法拉第定律两边再取旋度：
\[
\curl(\curl \vb{E})=-\pdv{}{t}(\curl \vb{B}).
\]
利用矢量恒等式
\[
\curl(\curl \vb{A})=\grad(\div \vb{A})-\nabla^2\vb{A},
\]
得
\[
\grad(\div \vb{E})-\nabla^2\vb{E}=-\pdv{}{t}(\curl \vb{B}).
\]
因为真空中 $\div \vb{E}=0$，所以上式左边化为
\[
-\nabla^2\vb{E}=-\pdv{}{t}(\curl \vb{B}).
\]
再代入安培--麦克斯韦方程
\[
\curl \vb{B}=\mu_0\epsilon_0\pdv{\vb{E}}{t},
\]
得
\[
-\nabla^2\vb{E}=-\mu_0\epsilon_0\pdv[2]{\vb{E}}{t}.
\]
于是
\[
\nabla^2\vb{E}-\mu_0\epsilon_0\pdv[2]{\vb{E}}{t}=0.
\]
同理，对安培--麦克斯韦方程取旋度：
\[
\curl(\curl \vb{B})=\mu_0\epsilon_0\pdv{}{t}(\curl \vb{E}).
\]
利用同一恒等式并用 $\div \vb{B}=0$，得到
\[
-\nabla^2\vb{B}=\mu_0\epsilon_0\pdv{}{t}\qty(-\pdv{\vb{B}}{t})
=-\mu_0\epsilon_0\pdv[2]{\vb{B}}{t}.
\]
因此
\[
\nabla^2\vb{B}-\mu_0\epsilon_0\pdv[2]{\vb{B}}{t}=0.
\]
这表明电磁场在真空中以速度
\[
\upsilon=\frac{1}{\sqrt{\mu_0\epsilon_0}}=c
\]
传播，即电磁波。
\end{solution}
\end{problem}

\begin{problem}
设球对称标量场满足三维波动方程
\[
\nabla^2\psi-\frac{1}{c^2}\pdv[2]{\psi}{t}=0,
\]
且 $\psi=\psi(r,t)$。证明球面波通解可写成
\[
\psi(r,t)=\frac{1}{r}f(r-ct)+\frac{1}{r}g(r+ct).
\]
并解释其物理意义。
\begin{solution}
在球对称条件下，拉普拉斯算符化为
\[
\nabla^2\psi=\frac{1}{r^2}\pdv{}{r}\qty(r^2\pdv{\psi}{r}).
\]
令
\[
u(r,t)=r\psi(r,t).
\]
则
\[
\psi=\frac{u}{r},
\qquad
\pdv{\psi}{r}=\frac{1}{r}\pdv{u}{r}-\frac{u}{r^2}.
\]
于是
\[
r^2\pdv{\psi}{r}=r\pdv{u}{r}-u,
\]
再对 $r$ 求导：
\[
\pdv{}{r}\qty(r^2\pdv{\psi}{r})=r\pdv[2]{u}{r}.
\]
所以
\[
\nabla^2\psi=\frac{1}{r}\pdv[2]{u}{r}.
\]
另一方面
\[
\pdv[2]{\psi}{t}=\frac{1}{r}\pdv[2]{u}{t}.
\]
代入原波动方程得
\[
\frac{1}{r}\qty(\pdv[2]{u}{r}-\frac{1}{c^2}\pdv[2]{u}{t})=0.
\]
故 $u(r,t)$ 满足一维波动方程：
\[
\pdv[2]{u}{r}-\frac{1}{c^2}\pdv[2]{u}{t}=0.
\]
其达朗贝尔通解为
\[
u(r,t)=f(r-ct)+g(r+ct).
\]
因此
\[
\psi(r,t)=\frac{1}{r}f(r-ct)+\frac{1}{r}g(r+ct).
\]
其中第一项表示向外传播的球面波，第二项表示向内传播的球面波。前面的 $1/r$ 因子说明波在传播时振幅按距离反比衰减，这是因为能量分布在面积随 $r^2$ 增大的球面上。
\end{solution}
\end{problem}
